\question[4]
Erläutere, was es bedeutet, dass eine Datenbank in der 1. Normalform vorliegt.
Gib ein Beispiel an für eine Datenbank, die nicht in der 1. Normalform vorliegt, 
und erkläre, wie man sie in die 1. Normalform überführen kann.


\begin{solutionbox}{10cm}   
Eine Datenbank befindet sich in der 1. Normalform, wenn alle Attribute nur atomare Werte enthalten.
Das bedeutet: In einer Zelle steht genau ein Wert, keine Listen, keine mehrfachen Werte und keine zusammengesetzten Felder.

\vspace{0.2cm}

\textbf{Beispiel für keine 1. Normalform:}

\begin{center}
\begin{tabular}{|l|l|}
\hline
Name & Telefonnummern \\
\hline
Anna Meier & 030-123456, 030-654321 \\
Max Schulz & 040-987654 \\
\hline
\end{tabular}
\end{center}

Hier ist das Attribut \textit{Telefonnummern} nicht atomar, weil mehrere Werte in einer Zelle stehen.

\vspace{0.2cm}

\textbf{Überführung in die 1. Normalform:}

\begin{center}
\begin{tabular}{|l|l|}
\hline
Name & Telefonnummer \\
\hline
Anna Meier & 030-123456 \\
Anna Meier & 030-654321 \\
Max Schulz & 040-987654 \\
\hline
\end{tabular}
\end{center}

Nun hat jede Zeile genau einen Telefonwert, also liegt die Tabelle in der 1. Normalform vor.

\end{solutionbox}